\section{Image analysis pipeline}
\section{Introduction}
	
	The petrographic polarising microscope is a foundational tool in geoscience research to answer first-order questions such as rock micro-structure, fabric, and mineral assemblage at multiple observation scales. The arrival of fast and reliable optical slide scanners for biomedical imaging has motivated their re-implementation as polarising microscopes for imaging rock thin sections. Investigators are now demanding sharing their slide data around the world via virtual microscopes, scaling up image analysis to hundreds of thin sections, and integrating optical imagery with other modalities, especially chemical maps.
	
	The literature shows image analysis of rock thin sections studying small fields of view and targeting mineral grains with one microscopy technique is a challenging task. Large-scale and correlative (more than one technique) microscopy requires a step evolution of image analysis software and algorithms that can cope with image pyramids for registration, segmentation, classification, and image representation of pixels and objects that simplify observation of the underlying data. The average user will need centralised orchestration of data management, image processing, and image analysis for trialling and interacting with open-source code and data that might not be locally available (the cloud).
	
	Segmentation has been demonstrated to work with QuPath (Bankhead et al., 2017) with extensive use of the Pixel Classifier following the elaboration of multi-channel images with Image Combiner Warpy. The segmentation of integrated reflected light (RL), PPL-max, and XPL-max only has achieved results comparable to SEM-based Automated mineralogy systems. Optical microscopy reaches good results as a standalone technique when the sample has contrasting pleochroism and interference colours that are ~invariant to crystal orientation.
	
	
	With this software you can add new dimensions to your analysis since polarised optical microscopy can rapidly map for:
	
	\begin{itemize}
		
		\item Mineral identification: distinctive colour (Acevedo Zamora & Kamber, 2023), isochemical phases
		\item Texture: grains (properties), boundaries, contacts, neighbouring relationships (Kamber et al., 2025), and rock fabric
		\item Fine micro-structure (e.g., diagenetic dykes), even if they are not shown in chemical maps (Acevedo Zamora et al., 2024)
		\item Accessory phases and micro-inclusions (at high magnification)
		\item Sample depth/volume (focusing as planes or extended depth of focus)
		\item Mineral optic-axis and/or slow-axis orientation (Acevedo Zamora et al., 2024)
		\item Sample preparation quality and planing the location of new micro-analytical experiments, e.g., XPL colours depend on thickness
	
	\end{itemize}
	
	
	\section{Installation}
		
	\section{Citations}
	
	The software depends on open-source as well (see above) and scientific citations/feedback. The following research papers already have contributed to its evolution:
	
	\begin{itemize}
		
		\item Cube Converter:	
		
			\begin{itemize}				
				\item Acevedo Zamora, M. A., \& Kamber, B. S. (2023). Petrographic Microscopy with Ray Tracing and Segmentation from Multi-Angle Polarisation Whole-Slide Images. Minerals, 13(2), 156. https://doi.org/10.3390/min13020156
				
				\item Acevedo Zamora, M. (2024). Petrographic microscopy of geologic textural patterns and element-mineral associations with novel image analysis methods [Thesis by publication, Queensland University of Technology]. Brisbane. https://eprints.qut.edu.au/248815/
			\end{itemize}
		
		\item Chemistry Simplifier:	
		
			\begin{itemize}
				\item Acevedo Zamora, M. A., Kamber, B. S., Jones, M. W. M., Schrank, C. E., Ryan, C. G., Howard, D. L., Paterson, D. J., Ubide, T., \& Murphy, D. T. (2024). Tracking element-mineral associations with unsupervised learning and dimensionality reduction in chemical and optical image stacks of thin sections. Chemical Geology, 650, 121997. https://doi.org/10.1016/j.chemgeo.2024.121997
				
				\item Acevedo Zamora, M. (2024). Petrographic microscopy of geologic textural patterns and element-mineral associations with novel image analysis methods [Thesis by publication, Queensland University of Technology]. Brisbane. https://eprints.qut.edu.au/248815/
				
				\item Ubide, T., Murphy, D. T., Emo, R. B., Jones, M. W. M., Acevedo Zamora, M. A., \& Kamber, B. S. (2025). Early pyroxene crystallisation deep below mid-ocean ridges. Earth and Planetary Science Letters, 663, 119423. https://doi.org/10.1016/j.epsl.2025.119423
						
			\end{itemize}
		
		\item Phase Interpreter:	
		
			\begin{itemize}
				\item Acevedo Zamora, M. A., \& Kamber, B. S. (2023). Petrographic Microscopy with Ray Tracing and Segmentation from Multi-Angle Polarisation Whole-Slide Images. Minerals, 13(2), 156. https://doi.org/10.3390/min13020156
				
				\item Acevedo Zamora, M. A., Kamber, B. S., Jones, M. W. M., Schrank, C. E., Ryan, C. G., Howard, D. L., Paterson, D. J., Ubide, T., \& Murphy, D. T. (2024). Tracking element-mineral associations with unsupervised learning and dimensionality reduction in chemical and optical image stacks of thin sections. Chemical Geology, 650, 121997. https://doi.org/10.1016/j.chemgeo.2024.121997
				
				\item Acevedo Zamora, M. (2024). Petrographic microscopy of geologic textural patterns and element-mineral associations with novel image analysis methods [Thesis by publication, Queensland University of Technology]. Brisbane. https://eprints.qut.edu.au/248815/
				
				\item Ubide, T., Murphy, D. T., Emo, R. B., Jones, M. W. M., Acevedo Zamora, M. A., \& Kamber, B. S. (2025). Early pyroxene crystallisation deep below mid-ocean ridges. Earth and Planetary Science Letters, 663, 119423. https://doi.org/10.1016/j.epsl.2025.119423
				
				\item Kamber, B. S., Acevedo Zamora, M. A., Rodrigues, R. F., Li, M., Yaxley, G. M., \& Ng, M. (2025). Exploring High PT Experimental Charges Through the Lens of Phase Maps. Minerals, 15(4), 355. https://doi.org/10.3390/min15040355
				
				\item Rodrigues, R. F., Yaxley, G. M., \& Kamber, B. S. (2025). Phase relations and solidus temperature of garnet lherzolite at 5 GPa revisited. Contributions to Mineralogy and Petrology, 180(9), 57. https://doi.org/10.1007/s00410-025-02250-4			
				
			\end{itemize}	
		
	\end{itemize}
	

	% --- REPOSITORY 1 ---